\begin{statementbox}
\paragraph{F001 导数怎样判断函数是增还是减}
\begin{center}\begin{center}
\begin{tikzpicture}
\begin{axis}[width=11.6cm,height=6.3cm,axis lines=middle,
 xmin=-1.5,xmax=1.65,ymin=-2.5,ymax=2.7,
 xlabel={$x$},ylabel={$y$},ticks=none,clip=false,
 title={同一条曲线上：正切线斜率与高度增加},title style={font=\small}]
\addplot[blue,very thick,domain=-1.35:1.35,samples=100]{x+x^3/3};
\addplot[orange!85!black,thick,dashed,domain=-1.3:-0.3]{1.64*(x+0.8)-0.9706667};
\addplot[orange!85!black,thick,dashed,domain=0.3:1.3]{1.64*(x-0.8)+0.9706667};
\addplot[only marks,blue] coordinates {(-0.8,-0.9706667) (0.8,0.9706667)};
\node[font=\small,anchor=north west] at (axis cs:-1.2,-1.65) {$f'(-0.8)>0$};
\node[font=\small,anchor=south west] at (axis cs:0.1,2) {$f'(0.8)>0$};
\draw[densely dotted] (axis cs:-0.8,0)--(axis cs:-0.8,-0.9706667);
\draw[densely dotted] (axis cs:0.8,0)--(axis cs:0.8,0.9706667);
\end{axis}
\end{tikzpicture}
\end{center}
\end{center}
在同一段区间内，内部可导、所含端点连续时：
\[
 f'(x)>0\ \Longrightarrow\ \text{严格递增},\qquad
 f'(x)<0\ \Longrightarrow\ \text{严格递减}.
\]
若只知道 $f'(x)\ge0$，先得到“不下降”；还要排除整段水平线，才能进一步得到严格递增。

\end{statementbox}
\clearpage
\begin{statementbox}

设函数 $f(x)$ 在一段连续的取值范围 $I$ 上有定义：在这段范围内没有断开，中间的点都能求导；如果包含左右端点，还要求函数在端点处连续。下面说的导数条件，都指这段范围内部的点。

\textbf{一、最常用的四条规则}
\begin{enumerate}
  \item 若 $f'(x)>0$ 始终成立，那么 $x$ 从小到大时，$f(x)$\textbf{严格递增}（一直往上走）。
  \item 若 $f'(x)<0$ 始终成立，那么 $f(x)$\textbf{严格递减}（一直往下走）。
  \item 若 $f'(x)\ge 0$ 始终成立，那么 $f(x)$\textbf{不会下降}，也就是函数值可以升高，也可以暂时保持不变。
  \item 若 $f'(x)\le 0$ 始终成立，那么 $f(x)$\textbf{不会上升}，也就是函数值可以降低，也可以暂时保持不变。
\end{enumerate}

\textbf{二、``不会下降''和``严格递增''有什么区别？}
任意取 $x_1<x_2$：
\[
\begin{aligned}
\text{严格递增}&:\quad f(x_1)<f(x_2),\\
\text{不会下降}&:\quad f(x_1)\le f(x_2).
\end{aligned}
\]
一个区别：\textbf{严格递增不允许两个不同位置的函数值相等；不会下降允许出现一整段水平线。}

\textbf{三、最容易忽略的重要补充}
\begin{enumerate}
  \item $f'(x)>0$ 是严格递增的\textbf{好用条件}，但不是必须处处满足的条件。例如 $f(x)=x^3$ 严格递增，却有 $f'(0)=0$。
  \item 若 $f'(x)\ge0$，并且\textbf{没有任何一小段区间让函数值始终不变}，那么 $f(x)$ 仍严格递增。
  \item 特别好用的判断：若 $f'(x)\ge0$，而且导数等于 $0$ 的位置\textbf{只有有限个}，那么 $f(x)$ 一定严格递增。
\end{enumerate}
\textbf{递减时同样成立：}把上面的 $\ge0$ 改为 $\le0$，把``严格递增''改为``严格递减''。
\end{statementbox}
